\documentclass[11pt,a4paper]{article}

% ── packages ──────────────────────────────────────────────
\usepackage[margin=1in]{geometry}
\usepackage{amsmath,amssymb,amsthm}
\usepackage{mathtools}
\usepackage{bm}
\usepackage{booktabs}
\usepackage{enumitem}
\usepackage{hyperref}
\usepackage[ruled,vlined]{algorithm2e}
\usepackage{natbib}
\usepackage{xcolor}
\hypersetup{colorlinks=true,linkcolor=blue!60!black,citecolor=blue!60!black,urlcolor=blue!50!black}

\newcommand{\R}{\mathbb{R}}
\newcommand{\bx}{\bm{x}}
\newcommand{\bs}{\bm{s}}
\newcommand{\bmu}{\bm{\mu}}
\newcommand{\bSigma}{\bm{\Sigma}}
\DeclareMathOperator{\diag}{diag}

\title{\textbf{Data-Generation Schemes for Bivariate Spatial Simulation}\\[4pt]
       \large Non-Stationary and Non-Gaussian Settings}
\author{}
\date{}

\begin{document}
\maketitle
\thispagestyle{empty}

%=====================================================================
\section{Common Setup}
%=====================================================================

Both simulation schemes share the following spatial domain and sampling design.

\begin{enumerate}[leftmargin=*]
  \item \textbf{Regular grid.}
    Define two marginal grids $x_k = k/79$, $k=0,1,\dots,79$ and
    $y_l = l/79$, $l=0,1,\dots,79$.
    The full grid is $\mathcal{G} = \{(x_k,y_l)\} \subset [0,1]^2$ with
    $|\mathcal{G}|=6400$.

  \item \textbf{Sub-sampling.}
    From $\mathcal{G}$ we draw $N=1200$ locations uniformly at random
    without replacement, yielding $\{\bs_i\}_{i=1}^{N}$.

  \item \textbf{Train / test split.}
    For every Monte-Carlo replicate, we randomly assign $N_{\text{train}}=1080$
    points to the training set and $N_{\text{test}}=120$ to the test set.

  \item \textbf{Number of replicates.}
    All results are averaged over $B=50$ independent simulation replicates.
\end{enumerate}

%=====================================================================
\section{Scheme~1: Non-Stationary Bivariate Field}
\label{sec:ns}
%=====================================================================

\subsection{Multi-Resolution Wendland Basis}

We construct a multi-resolution set of compactly-supported basis functions.
Let the number of knots at resolution $r$ be
$K_r \in \{4,\,9,\,25\}$ (i.e.\ $2^2,\,3^2,\,5^2$).
For resolution $r$, define a one-dimensional knot sequence
\[
  t_j^{(r)} = \frac{j}{\sqrt{K_r}-1}, \qquad j=0,1,\dots,\sqrt{K_r}-1,
\]
and form the two-dimensional knot set
$\mathcal{D}_r = \{(t_j^{(r)},\,t_k^{(r)})\}$.

For each knot $\bm{d}_m \in \mathcal{D}_r$ the Wendland $C^6$ radial basis function
evaluated at location $\bs$ is
\begin{equation}\label{eq:wendland}
  \phi_m(\bs) =
  \begin{cases}
    \dfrac{(1-\|\bs-\bm{d}_m\|)^{6}\,
          \bigl(35\|\bs-\bm{d}_m\|^{2}+18\|\bs-\bm{d}_m\|+3\bigr)}{3},
      & 0 \le \|\bs-\bm{d}_m\| \le 1,\\[6pt]
    0, & \text{otherwise}.
  \end{cases}
\end{equation}

Stacking all resolutions gives the $N\times P$ basis matrix
$\bm{\Phi}=[\phi_m(\bs_i)]$, where $P = 4+9+25 = 38$.

\subsection{Bivariate Response Generation}

Partition the columns of $\bm{\Phi}$ into even-indexed columns
$\bm{\Phi}_{\text{even}}$ and odd-indexed columns $\bm{\Phi}_{\text{odd}}$.
At each replicate, draw six coefficients independently and uniformly:
\[
  a,\,b,\,c,\,d,\,e,\,f \;\sim\; \text{Uniform}\!\bigl(\{-2.5,\,-2.449,\dots,2.5\}\bigr).
\]

The two response variables are
\begin{align}
  Z_1(\bs_i) &= \sum_m \Bigl[
        a\,\bigl(\bm{\Phi}_{\text{even}}\bigr)_{im}^{3/2}
      + c\,\bigl(\bm{\Phi}_{\text{odd}}\bigr)_{im}
      - b\,\sqrt{\bigl(\bm{\Phi}_{\text{even}}\bigr)_{im}
                 \bigl(\bm{\Phi}_{\text{odd}}\bigr)_{im}}
    \Bigr], \label{eq:ns_var1}\\[4pt]
  Z_2(\bs_i) &= \sum_m \Bigl[
        d\,\bigl(\bm{\Phi}_{\text{even}}\bigr)_{im}
      - f\,\bigl(\bm{\Phi}_{\text{odd}}\bigr)_{im}^{3/2}
    \Bigr]. \label{eq:ns_var2}
\end{align}

\noindent
Because the coefficients and nonlinear transforms change across replicates,
the resulting fields are non-stationary (the local mean and variance depend
on the spatial location~$\bs$).

%=====================================================================
\section{Scheme~2: Non-Gaussian Bivariate Field with Covariates}
\label{sec:ng}
%=====================================================================

\subsection{Bivariate Mat\'ern Cross-Covariance}

Let $\|\bs_i-\bs_j\|=h$.  The Mat\'ern correlation function is
\begin{equation}\label{eq:matern}
  \mathcal{M}(h;\,\alpha,\,\nu)
  = \frac{2^{1-\nu}}{\Gamma(\nu)}
    \Bigl(\frac{h}{\sqrt{\alpha}}\Bigr)^{\!\nu}
    K_{\nu}\!\Bigl(\frac{h}{\sqrt{\alpha}}\Bigr),
  \qquad h>0,
\end{equation}
with $\mathcal{M}(0;\alpha,\nu)=1$.
Here $K_\nu$ is the modified Bessel function of the second kind.

The $2N\times 2N$ block covariance matrix for the bivariate field is
\[
  \bSigma_{\text{biv}} =
  \begin{pmatrix}
    \bSigma_{11} & \bSigma_{12}\\
    \bSigma_{12}^\top & \bSigma_{22}
  \end{pmatrix},
\]
where
\begin{align}
  (\bSigma_{11})_{ij} &= \sigma_{11}\;\mathcal{M}(h;\,\alpha_{11},\,\nu_{11}), \\
  (\bSigma_{22})_{ij} &= \sigma_{22}\;\mathcal{M}(h;\,\alpha_{22},\,\nu_{22}), \\
  (\bSigma_{12})_{ij} &= \kappa\;\mathcal{M}(h;\,\alpha_{12},\,\nu_{12}).
\end{align}
The cross-covariance scaling constant is
\[
  \kappa
  = \frac{\sqrt{\sigma_{11}\sigma_{22}}\;R\;
          \Gamma(\nu_{12})}
         {\sqrt{\Gamma(\nu_{11})\,\Gamma(\nu_{22})}}
  \;\cdot\;
  \frac{\sqrt{\alpha_{11}^{\,\nu_{11}}\;\alpha_{22}^{\,\nu_{22}}}}
       {\alpha_{12}^{\,\nu_{12}}},
\]
with the parameter values listed in Table~\ref{tab:params}.

\begin{table}[h]
\centering
\caption{Bivariate Mat\'ern parameters.}\label{tab:params}
\smallskip
\begin{tabular}{@{}lcccccc@{}}
\toprule
 & $\sigma$ & $\nu$ & $\alpha$ & $R$ \\
\midrule
Variable 1 & 0.7 & 0.3 & 0.05 & --- \\
Variable 2 & 0.8 & 0.6 & 0.10 & --- \\
Cross      & --- & 0.45 & 0.075 & 0.5 \\
\bottomrule
\end{tabular}
\end{table}

\subsection{Spatial Covariates}

Five covariates $x_1,\dots,x_5$ are generated as independent realisations of a
Gaussian process with constant mean
$\mu_k \sim \text{Uniform}(\{-2.5,\dots,2.5\})$ and Mat\'ern covariance
$\sigma_c\,\mathcal{M}(h;\,0.1,\,0.5)$ with $\sigma_c=0.9$.

\subsection{Nonlinear Mean Function}

The shared (location-dependent) mean function is
\begin{equation}\label{eq:mean}
\begin{aligned}
  m(\bs) &= x_1^2 - x_2^2 + x_3^2 - x_4^2 - x_5^2
           + 2x_1 x_2 + 3x_2 x_3 - 2x_3 x_5 + 10 x_1 x_4 \\
         &\quad + \sin(x_1)\,x_2\,x_3
           + \cos(x_2)\,x_3\,x_5
           + x_1 x_2 x_4 x_5.
\end{aligned}
\end{equation}

\subsection{Tukey $g$--$h$ Transformation}

Draw a single bivariate Gaussian realisation
\[
  \begin{pmatrix} W_1 \\ W_2 \end{pmatrix}
  \sim \mathcal{N}\!\Bigl(\bm{0},\;\bSigma_{\text{biv}}\Bigr).
\]

The Tukey $g$--$h$ transform of a standard Gaussian variable $W$ with
parameters $(g,h)$ is
\begin{equation}\label{eq:tukey}
  T_{g,h}(W) = \frac{e^{gW}-1}{g}\;\exp\!\Bigl(\frac{h\,W^2}{2}\Bigr).
\end{equation}

The observed responses are
\begin{align}
  Z_1(\bs_i) &= T_{0.8,\;0.5}(W_{1i}) + m(\bs_i), \\
  Z_2(\bs_i) &= T_{-0.8,\;0.5}(W_{2i}) + m(\bs_i).
\end{align}

\noindent
The positive $g$ for $Z_1$ induces right skewness; the negative $g$ for~$Z_2$
induces left skewness.  The $h=0.5$ parameter produces heavy tails in both
margins, making the joint distribution substantially non-Gaussian.

%=====================================================================
\section{Estimation and Evaluation}
%=====================================================================

\subsection{Methods}

Given training data $\{(\bs_i,\bm{Z}_i)\}_{i=1}^{1080}$, we predict
$\hat{\bm{Z}}_j$ at each test location $\bs_j$ ($j=1,\dots,120$) using:

\begin{enumerate}[leftmargin=*]
  \item \textbf{Conditional GAN (cGAN).}
    A generator $G_\theta(\bs,\bm{z})$ maps spatial coordinates (and
    covariates, if available) plus noise $\bm{z}\sim\mathcal{N}(\bm{0},\bm{I}_{32})$
    to $\hat{\bm{Z}}\in\R^2$.  A discriminator $D_\psi$ distinguishes real
    from generated pairs.  At prediction time we draw $S=200$ samples per
    test location and report the empirical mean and standard deviation.

  \item \textbf{$K$-Nearest Neighbours (KNN).}
    Distance-weighted regression with $K=10$ in the (standardised)
    feature space.

  \item \textbf{Random Forest (RF).}
    An ensemble of 300 regression trees (maximum depth 15), fitted
    independently for each response variable.  Prediction uncertainty is
    obtained from across-tree variance.
\end{enumerate}

\subsection{Metrics}

\begin{enumerate}[leftmargin=*]
  \item \textbf{Mean Squared Error (MSE).}
    \[
      \text{MSE} = \frac{1}{N_{\text{test}}}
        \sum_{j=1}^{N_{\text{test}}} \|\bm{Z}_j - \hat{\bm{Z}}_j\|^2.
    \]

  \item \textbf{Mean Absolute Error (MAE).}
    \[
      \text{MAE} = \frac{1}{N_{\text{test}}}
        \sum_{j=1}^{N_{\text{test}}} \|\bm{Z}_j - \hat{\bm{Z}}_j\|_1.
    \]

  \item \textbf{Mahalanobis Distance (MD).}
    Let $\hat{\bSigma}$ be the empirical covariance of the residuals.  Then
    \[
      \text{MD} = \frac{1}{N_{\text{test}}}
        \sum_{j=1}^{N_{\text{test}}}
        \sqrt{(\bm{Z}_j-\hat{\bm{Z}}_j)^\top
              \hat{\bSigma}^{-1}
              (\bm{Z}_j-\hat{\bm{Z}}_j)}.
    \]

  \item \textbf{Continuous Ranked Probability Score (CRPS).}
    For a Gaussian predictive distribution
    $\hat{Z}_{jv}\sim\mathcal{N}(\hat\mu_{jv},\hat\sigma_{jv}^2)$:
    \[
      \text{CRPS}(\hat\mu,\hat\sigma,z)
      = \hat\sigma\!\left[
          \zeta\,\Phi(\zeta) + \varphi(\zeta) - \frac{1}{\sqrt\pi}
        \right],
      \qquad \zeta = \frac{z-\hat\mu}{\hat\sigma},
    \]
    where $\Phi$ and $\varphi$ are the standard normal CDF and PDF.

  \item \textbf{95\% Coverage (COV95).}
    Fraction of test observations falling within
    $[\hat\mu - 1.96\hat\sigma,\;\hat\mu + 1.96\hat\sigma]$.
    Ideal value: 0.95.
\end{enumerate}

%=====================================================================
\section{Summary}
%=====================================================================

Table~\ref{tab:overview} summarises the two simulation schemes.

\begin{table}[h]
\centering
\caption{Overview of the two data-generation schemes.}\label{tab:overview}
\smallskip
\begin{tabular}{@{}lcc@{}}
\toprule
Feature & Scheme 1 (Non-Stationary) & Scheme 2 (Non-Gaussian) \\
\midrule
Spatial domain     & $[0,1]^2$       & $[0,1]^2$ \\
Grid size          & $80\times80$    & $80\times80$ \\
Sample size $N$    & 1200            & 1200 \\
Training / test    & 1080 / 120      & 1080 / 120 \\
Basis / covariance & Wendland $C^6$  & Bivariate Mat\'ern \\
Non-Gaussianity    & ---             & Tukey $g$--$h$ \\
Covariates         & None            & 5 spatial GPs \\
Mean function      & Nonlinear basis & Nonlinear (Eq.~\ref{eq:mean}) \\
Replicates $B$     & 50              & 50 \\
\bottomrule
\end{tabular}
\end{table}

\end{document}
